\chapter{Modelisation Merise}

\section{Introduction}

La methode Merise permet de modeliser les donnees du systeme a travers trois niveaux
d'abstraction : conceptuel (MCD), logique (MLD) et physique (MPD). Notre systeme MedSys
etant base sur une architecture microservices, nous presentons la modelisation pour chaque
microservice separement.

% ============================================
\section{Modele Conceptuel de Donnees (MCD)}

\subsection{MCD -- Microservice ms-auth}

\begin{figure}[H]
  \centering
  \includegraphics[width=0.6\textwidth]{figures/mcd_auth}
  \caption{MCD du microservice ms-auth}
  \label{fig:mcd_auth}
\end{figure}

Le microservice d'authentification repose sur une entite centrale
\texttt{USER\_ACCOUNT} qui centralise les informations de connexion
et les references vers les entites metier des autres microservices
(\texttt{patientId} vers ms-patient, \texttt{personnelId} vers ms-personnel).

\subsection{MCD -- Microservice ms-patient-personnel}

\begin{figure}[H]
  \centering
  \includegraphics[width=\textwidth]{figures/mcd_patient}
  \caption{MCD du microservice ms-patient-personnel}
  \label{fig:mcd_patient}
\end{figure}

Le dossier medical (\texttt{DOSSIER\_MEDICAL}) constitue l'entite centrale du
microservice patient. Il agrege l'ensemble des informations cliniques a travers
plusieurs entites specialisees : consultations, antecedents, ordonnances, analyses
de laboratoire, radiologies, hospitalisations, certificats medicaux, documents
uploades et messages.

Les cardinalites principales sont :
\begin{itemize}
  \item \textbf{PATIENT -- DOSSIER\_MEDICAL} : (1,1) -- (1,1) -- chaque patient possede exactement un dossier medical
  \item \textbf{DOSSIER\_MEDICAL -- CONSULTATION} : (1,1) -- (0,n) -- un dossier peut contenir zero ou plusieurs consultations
  \item \textbf{ORDONNANCE -- LIGNE\_ORDONNANCE} : (1,1) -- (1,n) -- une ordonnance contient au moins une ligne
  \item \textbf{MEDECIN -- SPECIALITE} : (0,n) -- (1,1) -- un medecin appartient a une specialite
\end{itemize}

% ============================================
\section{Modele Logique de Donnees (MLD)}

Le MLD traduit le MCD en tables relationnelles en introduisant les cles primaires,
cles etrangeres et contraintes d'integrite.

\subsection{MLD -- ms-auth}

\begin{verbatim}
USER_ACCOUNT (#id, email*, password, role, nom, prenom, cin,
              enabled, reset_token, reset_token_expiry,
              patient_id, personnel_id, created_at, updated_at)
\end{verbatim}

\subsection{MLD -- ms-patient-personnel}

\begin{verbatim}
SPECIALITE      (#id, nom*, code*)
SERVICE         (#id, nom, code*, description, localisation, capacite_lits)

MEDECIN         (#id, nom, prenom, matricule,
                 #specialite_id => SPECIALITE,
                 #service_id => SERVICE,
                 derniere_synchronisation)

PATIENT         (#id, nom, prenom, cin*, date_naissance, sexe,
                 groupe_sanguin, telephone, email, adresse, ville,
                 mutuelle, numero_cnss,
                 #dossier_medical_id => DOSSIER_MEDICAL,
                 created_at, updated_at)

DOSSIER_MEDICAL (#id, numero_dossier*, date_creation)

CONSULTATION    (#id, #dossier_id => DOSSIER_MEDICAL,
                 #medecin_id => MEDECIN,
                 date_consultation, motif, diagnostic,
                 observations, traitement, poids, taille,
                 tension_systolique, tension_diastolique, temperature)

ANTECEDENT      (#id, #dossier_id => DOSSIER_MEDICAL,
                 type_antecedent, description, date_diagnostic,
                 severite, actif, source)

ORDONNANCE      (#id, #dossier_id => DOSSIER_MEDICAL,
                 #medecin_id => MEDECIN,
                 date_ordonnance, type_ordonnance, instructions,
                 est_renouvele, date_expiration)

LIGNE_ORDONNANCE(#id, #ordonnance_id => ORDONNANCE,
                 medicament, dosage, posologie,
                 duree_jours, quantite, instructions)

ANALYSE_LABORATOIRE(#id, #dossier_id => DOSSIER_MEDICAL,
                    date_analyse, date_resultat, type_analyse,
                    resultats, valeur_reference, statut, severite,
                    laboratoire, prescripteur, chemin_fichier)

RADIOLOGIE      (#id, #dossier_id => DOSSIER_MEDICAL,
                 date_examen, type_examen, description,
                 conclusion, prescripteur, radiologue)

HOSPITALISATION (#id, #dossier_id => DOSSIER_MEDICAL,
                 #service_id => SERVICE,
                 #medecin_responsable_id => MEDECIN,
                 date_entree, date_sortie, motif, diagnostic,
                 numero_chambre, compte_rendu, actif)

CERTIFICAT_MEDICAL(#id, #dossier_id => DOSSIER_MEDICAL,
                   #medecin_id => MEDECIN,
                   type_certificat, date_emission,
                   contenu, duree_jours)

DOCUMENT_PATIENT(#id, #dossier_id => DOSSIER_MEDICAL,
                 type_document, nom_fichier_original,
                 nom_fichier_stocke, chemin_fichier,
                 description, taille_fichier, content_type,
                 date_upload)

MESSAGE_PATIENT (#id, #dossier_id => DOSSIER_MEDICAL,
                 contenu, expediteur, lu,
                 medecin_id, medecin_nom, date_envoi)
\end{verbatim}

\textbf{Conventions :} \texttt{\#} = cle primaire, \texttt{*} = contrainte d'unicite, \texttt{=>} = cle etrangere.

% ============================================
\section{Modele Physique de Donnees (MPD)}

Le MPD correspond aux scripts SQL de creation des tables. Le script complet est
disponible dans le fichier \texttt{docs/sql/mpd.sql}.

\subsection{Tables principales -- ms\_auth\_db}

\begin{lstlisting}[caption=Table user\_accounts]
CREATE TABLE user_accounts (
    id        BIGINT AUTO_INCREMENT PRIMARY KEY,
    email     VARCHAR(255) NOT NULL UNIQUE,
    password  VARCHAR(255) NOT NULL,
    role      ENUM('PATIENT','MEDECIN','ADMIN',
                   'PERSONNEL','DIRECTEUR') NOT NULL,
    nom       VARCHAR(255) NOT NULL,
    prenom    VARCHAR(255) NOT NULL,
    cin       VARCHAR(50),
    enabled   BOOLEAN DEFAULT TRUE,
    ...
);
\end{lstlisting}

\subsection{Tables principales -- ms\_patient\_db}

\begin{lstlisting}[caption=Table patients]
CREATE TABLE patients (
    id                 BIGINT AUTO_INCREMENT PRIMARY KEY,
    cin                VARCHAR(50) NOT NULL UNIQUE,
    nom                VARCHAR(255) NOT NULL,
    prenom             VARCHAR(255) NOT NULL,
    date_naissance     DATE NOT NULL,
    dossier_medical_id BIGINT UNIQUE,
    FOREIGN KEY (dossier_medical_id)
        REFERENCES dossiers_medicaux(id)
);
\end{lstlisting}

\begin{lstlisting}[caption=Table dossiers\_medicaux]
CREATE TABLE dossiers_medicaux (
    id              BIGINT AUTO_INCREMENT PRIMARY KEY,
    numero_dossier  VARCHAR(50) NOT NULL UNIQUE,
    date_creation   DATETIME DEFAULT CURRENT_TIMESTAMP
);
\end{lstlisting}

\begin{lstlisting}[caption=Table consultations]
CREATE TABLE consultations (
    id                  BIGINT AUTO_INCREMENT PRIMARY KEY,
    dossier_id          BIGINT NOT NULL,
    medecin_id          BIGINT,
    date_consultation   DATETIME NOT NULL,
    motif               VARCHAR(500),
    diagnostic          TEXT,
    FOREIGN KEY (dossier_id)
        REFERENCES dossiers_medicaux(id),
    FOREIGN KEY (medecin_id)
        REFERENCES medecins_ref(id)
);
\end{lstlisting}

\textit{Le script SQL complet avec toutes les tables est disponible dans \texttt{docs/sql/mpd.sql}.}
